\section{AWS}
\tocc
\subsection{AWSとは}
\begin{frame}[t]{\ftitle}
    \begin{block}{AWS}
        Amazon Web Service の略称．Amazonが提供するクラウドサービスで，ネットワークを経由して仮想コンピュータやストレージなどのサービスを提供している．\hfill\cite{2015amazon}
    \end{block}
    \begin{minipage}{\textwidth}
        \centering
        \begin{minipage}[t]{.32\textwidth}
            \centering
            \includegraphics[keepaspectratio,width=\textwidth]{Arch_Amazon-EC2_64@5x.png}\\
            {\tiny Amazon Elastic Compute Cloud (Amazon EC2)\cite{aws}}
        \end{minipage}
        \begin{minipage}[t]{.32\textwidth}
            \centering
            \includegraphics[keepaspectratio,width=\textwidth]{Arch_Amazon-Simple-Storage-Service_64@5x.png}
            {\tiny Amazon Simple Storage Service (Amazon S3)\cite{aws}}
        \end{minipage}
        \begin{minipage}[t]{.32\textwidth}
            \centering
            \includegraphics[keepaspectratio,width=\textwidth]{Arch_Amazon-RDS_64@5x.png}\\
            {\tiny Amazon Relational Database Service (Amazon RDS)\cite{aws}}
        \end{minipage}
    \end{minipage}
\end{frame}
\subsection{オンプレミスとAWS}
\begin{frame}[t]{\ftitle}
    リソースを効率よく運用できる．\hfill\cite{2015amazon}
    \begin{figure}
        \centering
        \begin{minipage}[b]{.48\textwidth}
            \centering
            \begin{tikzpicture}
                \coordinate (O) at (0,0);
                \coordinate (A) at (1,2);
                \coordinate (B) at (2,1.5);
                \coordinate (C) at (3,3);
                \coordinate (D) at (4,2);
                \draw[red,thick, name path=curve]plot[smooth, tension=0.7]coordinates {(O)(A)(B)(C)(D)};
                \draw[name path=line,thick,dashed](0,2.5)--(4,2.5);
                \draw[thick,dashed](0,2.5)--(4,2.5);
                \draw[thick,-latex](0,0)--(4,0);
                \draw[thick,-latex](0,0)--(0,4.5);
                \node[below] at ($(0,0)!0.5!(4,0)$){\scriptsize 時間};
                \node[left] at ($(0,0)!0.5!(0,4.5)$){\scriptsize\rotatebox{90}{リソース量}};
                \node[below,text=red]at(B){\scriptsize 実際のリソース};
                \node[above]at(C){\scriptsize 機会損失};
                \node[below right]at(0,2.5){\scriptsize 無駄になる};
                \begin{scope}[on background layer]
                    \clip plot[smooth, tension=0.7] coordinates {(O) (A) (B) (C) (D)} |- (0,4) -- cycle;
                    \fill[yellow, name intersections={of=curve and line}] (intersection-1) rectangle (0,0);
                    \fill[yellow, name intersections={of=curve and line}] (intersection-2) rectangle (4,0);
                \end{scope}
            \end{tikzpicture}\\
            \vspace{1em}
            オンプレミス
        \end{minipage}
        \begin{minipage}[b]{.48\textwidth}
            \centering
            \begin{tikzpicture}
                \coordinate (O) at (0,0);
                \coordinate (A) at (1,2);
                \coordinate (B) at (2,1.5);
                \coordinate (C) at (3,3);
                \coordinate (D) at (4,2);
                \newcommand{\avv}{0.2cm}
                \draw[red,thick]plot[smooth, tension=0.7]coordinates {(O)(A)(B)(C)(D)};
                \draw[thick,dashed]plot[smooth,tension=0.7] coordinates {($(O)+(0,\avv)$)($(A)+(0,\avv)$)($(B)+(0,\avv)$)($(C)+(0,\avv)$)($(D)+(0,\avv)$)};
                \draw[thick,-latex](0,0)--(4,0);
                \draw[thick,-latex](0,0)--(0,4.5);
                \node[below] at ($(0,0)!0.5!(4,0)$){\scriptsize 時間};
                \node[left] at ($(0,0)!0.5!(0,4.5)$){\scriptsize\rotatebox{90}{リソース量}};
                \node[above]at($(C)+(0,\avv)$){\scriptsize 用意したリソース};
                \node[below,text=red]at(B){\scriptsize 必要なリソース};
            \end{tikzpicture}\\
            \vspace{1em}
            AWSの利用
        \end{minipage}
    \end{figure}
\end{frame}
\subsection{レンタルサーバとAWS}
\begin{frame}[t]{\ftitle}
    \begin{block}{レンタルサーバ}
        ``1台のサーバを複数のユーザが共用で利用する形態''\hfill\cite{2015amazon}
    \end{block}
    \begin{table}
        \fontsize{10pt}{10pt}\selectfont
        \centering
        \renewcommand{\arraystretch}{1.3}
        \newcommand{\hobo}{ほぼ}
        \begin{tabularx}{\textwidth}{cRRR}
                             & \multicolumn{1}{c}{共用サーバ} & \multicolumn{1}{c}{専用サーバ} & \multicolumn{1}{c}{仮想専用サーバ}          \\
            \hline
            {\tiny 利用形態}     & 1台の物理サーバを分割して利用           & 1台の物理サーバを占有               & 1台の物理サーバ上に立てた仮想サーバを利用                \\
            \hline
            {\tiny コスト}      & 安い                        & 高い                        & \textcolor{red}{中間}                  \\
            \hline
            {\tiny カスタマイズ}   & \hobo 不可能                 & 可能                        & \textcolor{red}{可能}                  \\
            \hline
            {\tiny セキュリティ}   & コントロール\textbf{不可}         & 高くコントロール\textbf{可能}       & \textcolor{red}{高くコントロール\textbf{可能}} \\
            \hline
            {\tiny 他のユーザの影響} & 強く受ける                     & 受けない                      & 多少受ける                                \\
            \hline
        \end{tabularx}
    \end{table}
    \hfill 引用\cite{2015amazon}
\end{frame}
\subsection{AWS導入例}
\begin{frame}[t]{\ftitle}
    \begin{itemize}
        \item Cloud LaTeX（アカリク）
              \begin{itemize}
                  \item コンパイルなど（EC2）
                  \item プロジェクト，テンプレート管理（S3，MySQL）
              \end{itemize}
        \item 株式会社任天堂
              \begin{itemize}
                  \item 『マリオカート ツアー』のDBに，「Amazon Aurora」（RDMS） を採用．
              \end{itemize}
        \item 株式会社東京海上日動火災保険
        \item JCB
        \item 盛岡市
        \item 浜松市
        \item SUBARU
    \end{itemize}
    \hfill\cite{導入事例}
\end{frame}